\documentclass{amsart}
\usepackage[utf8]{inputenc}
    
\usepackage{titlesec}
\titleformat{\section}{\normalfont\bf\Large}{\relax }{0em}{}
\titleformat{\subsection}{\normalfont\bf\large }{\relax}{0em}{}
\titleformat{\subsubsection}{\normalfont\bf}{\relax}{0em}{}

\usepackage{enumitem}
\usepackage{quoting}
\usepackage{array}
\usepackage{xcolor}
\long\def\gray#1{{\color{gray}{#1}}}\let\grey=\gray
\long\def\blue#1{{\color{blue}{#1}}}
\long\def\red#1{{\color{red}{#1}}}
\long\def\green#1{{\color{green}{#1}}}
\long\def\violet#1{{\color{violet}{#1}}}
\long\def\lila#1{{\color{lila}{#1}}}
\long\def\cyan#1{{\color{cyan}{#1}}}

\usepackage{wrapfig}
\usepackage{lpic}

\usepackage{manyfoot}
\SetFootnoteHook{\hspace{1.5em}}
\DeclareNewFootnote{A}
\SetFootnoteHook{\hspace{1.5em}}
\DeclareNewFootnote{E}[alph]
\usepackage[hyperfootnotes=false,
    colorlinks,
    linkcolor={blue!80!black},
    citecolor={blue!50!black},
    urlcolor={blue!80!black}]{hyperref}
\def\foothref#1#2{\href{#1}{ #2}\footnoteA{\url{#1}}}

\newcommand{\skipthis}[2][\dots]{\red{[#1]}\gray{#2}}
\long\def\skipthis#1{\relax}

\newtheorem*{theorem}{Theorem}
\newtheorem*{proposition}{Proposition}
\newtheorem*{question}{Question}
\newtheorem*{conjecture}{Conjecture}
\newtheorem*{definition}{Definition}

\let\dmo=\DeclareMathOperator

\let\tmp\epsilon
\let\epsilon\varepsilon
\let\varepsilon\tmp


\begin{document}
\centerline{\bf\Large From Topography to Geometry%
  \footnoteE{English translation by Rostislav Matveev of the article:
    Patrick Popescu-Pampu,  \textit{De la topographie \`a la
      g\'eom\'trie I et II}, Images des
    Math\'ematiques (2017);\\
    \url{https://doi.org/10.60868/7xv4-9g98} and
    \url{https://doi.org/10.60868/tzg9-3c39} (Ed.)}}%
\smallskip

\centerline{by Patrick Popescu-Pampu}
\bigskip

\centerline{\bf\large Contribution of Maps to Mathematical Thinking}
\medskip
% \begin{center}
  % \begin{lpic}[draft,b(12ex),t(1ex),clean]{Pix/ramsey(0.65)}
     % \lbl[t]{70,-5;\parbox{0.9\textwidth}{
        % \begin{center}
          % \sf Ordnung muss sein! (Order is inevitable.) This saying,
          % attributed to Theodore Motzkin, aptly summarizes Ramsey
          % theory, which aims to solve problems of the form: ``From what
          % size on can a structure no longer avoid having a certain
          % property?''%
          % \footnotemarkA{}
        % \end{center}
      % }}
  % \end{lpic}
% \end{center}











\end{document}
